% GENERATED FILE. Edit canonical lessons, then run npm run generate:books.
% canonical-source-hash: fnv1a64:d5848f7bf4a5124e
% canonical-lessons: LA-C405-abacus, LA-C405-calculus, LA-C405-enumero, LA-C405-dicto, LA-C405-pensum

\chapter{Counting and Dictation}
\label{ch:la-counting-and-dictation}
\hlchaptermodality{405}

\begin{chapteropening}
\textbf{By the end of this chapter:} \emph{I can talk about a counting board and pebbles for counting, listing things one by one, dictating to pupils and a set task.}

Complete the last lesson of chapter 405: I can talk about a counting board and pebbles for counting, listing things one by one, dictating to pupils and a set task.
\end{chapteropening}

\section[abacus, abacī]{abacus, abacī — a counting board; a sideboard}

\label{lesson:LA-C405-abacus}



\begin{quote}
Before the new one: say the Latin for geometry, then the Latin for a first principle.
\end{quote}

\subsection*{You'll want to know: abacus, abacī}
\textbf{abacus, abacī} — \textquotedblleft{}a counting board; a sideboard\textquotedblright{}.

\begin{quote}
Puer in abacō numerōs computat.
\end{quote}

\begin{quote}
\emph{The boy works out sums on the counting board.}
\end{quote}

\emph{abacus, abacī} m., borrowed from Greek \emph{abax}, a board. Pebbles moved along lines on the board. English \emph{abacus}.

\begin{cousinweb}
\emph{abacus} was borrowed from Greek \emph{abax, abakos} \textquotedblleft{}a board, a counting table\textquotedblright{} (EtymDB; Lewis \& Short). English took it over as \emph{abacus}.
\end{cousinweb}

\begin{grammarlens}[title={What you've built}]
One more word for telling, explaining and describing at B1.
\end{grammarlens}

\subsection*{Your turn}
\begin{itemize}
\raggedright
  \item \emph{Say it:} \emph{abacus, abacī}
  \item \emph{Say it:} \emph{abacus, abacī}, once more
  \item \emph{Say it:} a sentence of your own with \emph{abacus, abacī}
  \item \emph{Recall:} say the Latin for geometry, then the Latin for a first principle, then say \emph{abacus, abacī} again
\end{itemize}

\begin{tcolorbox}[breakable,colback=teal!4,colframe=teal!35!black,title={Before you move on}]
What does abacus, abacī mean? (\textquotedblleft{}A counting board; a sideboard\textquotedblright{}.) Say it once more.
\end{tcolorbox}
\section[calculus, calculī]{calculus, calculī — a pebble, a counter; a reckoning}

\label{lesson:LA-C405-calculus}



\begin{quote}
Before the new one: say the Latin for a first principle, then the Latin for a counting board.
\end{quote}

\subsection*{You'll want to know: calculus, calculī}
\textbf{calculus, calculī} — \textquotedblleft{}a pebble, a counter; a reckoning\textquotedblright{}.

\begin{quote}
Calculīs numerābat.
\end{quote}

\begin{quote}
\emph{He was counting with pebbles.}
\end{quote}

\emph{calculus, calculī} m., a little stone. \emph{Calculum pōnere}, to reckon. English \emph{calculate} and \emph{calculus}.

\begin{cousinweb}
\emph{calculus} is a small form of \emph{calx, calcis} \textquotedblleft{}a stone, lime\textquotedblright{} (Lewis \& Short): a little pebble used for counting. English \emph{calculate} comes from it.
\end{cousinweb}

\begin{grammarlens}[title={What you've built}]
One more word for telling, explaining and describing at B1.
\end{grammarlens}

\subsection*{Your turn}
\begin{itemize}
\raggedright
  \item \emph{Say it:} \emph{calculus, calculī}
  \item \emph{Say it:} \emph{calculus, calculī}, once more
  \item \emph{Say it:} a sentence of your own with \emph{calculus, calculī}
  \item \emph{Recall:} say the Latin for a first principle, then the Latin for a counting board, then say \emph{calculus, calculī} again
\end{itemize}

\begin{tcolorbox}[breakable,colback=teal!4,colframe=teal!35!black,title={Before you move on}]
What does calculus, calculī mean? (\textquotedblleft{}A pebble, a counter; a reckoning\textquotedblright{}.) Say it once more.
\end{tcolorbox}
\section[ēnumerō, ēnumerāre]{ēnumerō, ēnumerāre — to count out, to list}

\label{lesson:LA-C405-enumero}



\begin{quote}
Before the new one: say the Latin for a counting board, then the Latin for a pebble.
\end{quote}

\subsection*{You'll want to know: ēnumerō, ēnumerāre}
\textbf{ēnumerō, ēnumerāre} — \textquotedblleft{}to count out, to list\textquotedblright{}.

\begin{quote}
Omnia beneficia eius ēnumerāvit.
\end{quote}

\begin{quote}
\emph{He listed all his kindnesses.}
\end{quote}

\emph{ēnumerō, ēnumerāre, ēnumerāvī, ēnumerātum}, \emph{ē} + \emph{numerō}. English \emph{enumerate}.

\begin{cousinweb}
\emph{ēnumerō} is built from \emph{ē} \textquotedblleft{}out\textquotedblright{} and \emph{numerō} \textquotedblleft{}to count\textquotedblright{}, from \emph{numerus} \textquotedblleft{}a number\textquotedblright{} (Lewis \& Short). English \emph{enumerate} comes from it.
\end{cousinweb}

\begin{grammarlens}[title={What you've built}]
One more word for telling, explaining and describing at B1.
\end{grammarlens}

\subsection*{Your turn}
\begin{itemize}
\raggedright
  \item \emph{Say it:} \emph{ēnumerō, ēnumerāre}
  \item \emph{Say it:} \emph{ēnumerō, ēnumerāre}, once more
  \item \emph{Say it:} a sentence of your own with \emph{ēnumerō, ēnumerāre}
  \item \emph{Recall:} say the Latin for a counting board, then the Latin for a pebble, then say \emph{ēnumerō, ēnumerāre} again
\end{itemize}

\begin{tcolorbox}[breakable,colback=teal!4,colframe=teal!35!black,title={Before you move on}]
What does ēnumerō, ēnumerāre mean? (\textquotedblleft{}To count out, to list\textquotedblright{}.) Say it once more.
\end{tcolorbox}
\section[dictō, dictāre]{dictō, dictāre — to say repeatedly; to dictate}

\label{lesson:LA-C405-dicto}



\begin{quote}
Before the new one: say the Latin for a pebble, then the Latin for to count out.
\end{quote}

\subsection*{You'll want to know: dictō, dictāre}
\textbf{dictō, dictāre} — \textquotedblleft{}to say repeatedly; to dictate\textquotedblright{}.

\begin{quote}
Magister versūs discipulīs dictābat.
\end{quote}

\begin{quote}
\emph{The teacher used to dictate verses to the pupils.}
\end{quote}

\emph{dictō, dictāre, dictāvī, dictātum}, built on \emph{dīcō}: to say over and over, for someone to take down. English \emph{dictate}.

\begin{cousinweb}
\emph{dictō} is built on \emph{dīcō} \textquotedblleft{}to say\textquotedblright{} (EtymDB; Lewis \& Short), with the ending \emph{-tō} of a repeated act. English \emph{dictate} comes from it.
\end{cousinweb}

\begin{grammarlens}[title={What you've built}]
One more word for telling, explaining and describing at B1.
\end{grammarlens}

\subsection*{Your turn}
\begin{itemize}
\raggedright
  \item \emph{Say it:} \emph{dictō, dictāre}
  \item \emph{Say it:} \emph{dictō, dictāre}, once more
  \item \emph{Say it:} a sentence of your own with \emph{dictō, dictāre}
  \item \emph{Recall:} say the Latin for a pebble, then the Latin for to count out, then say \emph{dictō, dictāre} again
\end{itemize}

\begin{tcolorbox}[breakable,colback=teal!4,colframe=teal!35!black,title={Before you move on}]
What does dictō, dictāre mean? (\textquotedblleft{}To say repeatedly; to dictate\textquotedblright{}.) Say it once more.
\end{tcolorbox}
\section[pēnsum, pēnsī]{pēnsum, pēnsī — a task, a day's work}

\label{lesson:LA-C405-pensum}



\begin{quote}
Before the new one: say the Latin for to count out, then the Latin for to say repeatedly.
\end{quote}

\subsection*{You'll want to know: pēnsum, pēnsī}
\textbf{pēnsum, pēnsī} — \textquotedblleft{}a task, a day's work\textquotedblright{}.

\begin{quote}
Ancillae pēnsum suum perfēcērunt.
\end{quote}

\begin{quote}
\emph{The maids finished their day's work.}
\end{quote}

\emph{pēnsum, pēnsī} n., the participle of \emph{pendō}, to weigh: the wool weighed out to a spinner each day, then any set task.

\begin{cousinweb}
\emph{pēnsum} is the participle of \emph{pendō} \textquotedblleft{}to weigh out\textquotedblright{} (EtymDB; Lewis \& Short): the wool weighed out for a day's spinning.
\end{cousinweb}

\begin{grammarlens}[title={What you've built}]
One more word for telling, explaining and describing at B1.
\end{grammarlens}

\subsection*{Your turn}
\begin{itemize}
\raggedright
  \item \emph{Say it:} \emph{pēnsum, pēnsī}
  \item \emph{Say it:} \emph{pēnsum, pēnsī}, once more
  \item \emph{Say it:} a sentence of your own with \emph{pēnsum, pēnsī}
  \item \emph{Recall:} say the Latin for to count out, then the Latin for to say repeatedly, then say \emph{pēnsum, pēnsī} again
\end{itemize}

\begin{tcolorbox}[breakable,colback=teal!4,colframe=teal!35!black,title={Before you move on}]
What does pēnsum, pēnsī mean? (\textquotedblleft{}A task, a day's work\textquotedblright{}.) Say it once more.
\end{tcolorbox}
